\scanpage{28}
On a donc une \uline{extension}
\[
1\longrightarrow G'_S(S)\unsure{\longrightarrow|\mathbin{-}}\unsure{\mathcal E|E}\longrightarrow\Gamma\longrightarrow1\mathbin{\unsure{-|\longrightarrow}}\unsure{0|\circ|\ldotp}
\tag*{($\mathcal E$)}
\]

\deleted{Supposons que $\inserted{\pi}\pi\unsure{\in|\subset}\Gamma$ opère sur}

Les \uline{sections} de celle-ci sont les structures de torseur \unsure{sous|sur} $(S,\Gamma;G_S)$ sur $Q_S$ qui prolongent sa structure de Tor. sur $(S,G_S)$.

Si on a une telle section, on en déduit une opération de $\Gamma$ sur $G'_S(S)$ (qui en fait provient de son opération sur $G'_S$ lui-\unsure{=|même} !) et les autres sections correspondent aux 1-cocycles de $\Gamma$ à valeurs dans $G'_S(S)$. \deleted{Les}\inserted{L’ens $\Gamma(\mathcal E)$ des} sections \unsure{forment|forme} \deleted{\illegible{1}} un ens où $G'_S(S)$ opère (autom. intérieurs) et deux \deleted{\unsure{objets sont}} $\Gamma(\unsure{C|\mathcal E})$ sont \inserted{conjugués ssi le 1-cocycle faisant passer de l’un à l’autre est \unsure{``trivial''|réduit|cohom. à|coh. à} 1} Si on a un sous-groupe $\pi\subset\Gamma$ et une section $\alpha_0:\pi\longrightarrow\mathcal E$ partielle, on peut s’intéresser à l’ens des $g'\in G'_S(\unsure{S|1|L})$ qui la fixent, c’est \unsure{l’ens} des éléments de $G'_S(S)$ fixés par $\pi$ (y opérant comme de juste) $G'_S(S)^\pi\,(=\!G'_L(L))$. Les sections $\alpha$ qui prolongent $\alpha_0$ sont donc telles que $G'_S(S)^\pi$ y opère, Si on a
% archive sheet number: 27
